La distribución de la antigüedad laboral entre los clientes con y sin mora presenta un comportamiento muy similar. Ambos grupos muestran valores mínimos y máximos comparables, así como una dispersión semejante. La principal diferencia se observa en la concentración de los datos, donde los clientes con mora presentan una mayor densidad de observaciones aproximadamente entre los
6.5 y 22 años
de antigüedad laboral. Sin embargo, esta diferencia resulta poco pronunciada y no permite identificar un patrón claramente diferenciador entre ambos grupos.
A partir de este análisis descriptivo, la antigüedad laboral por sí sola no parece constituir un factor suficientemente explicativo del comportamiento de la mora. En consecuencia, se recomienda complementar este estudio incorporando variables adicionales, como el nivel de ingresos, el monto del préstamo, la relación cuota-ingreso y otras características del cliente, así como aplicar técnicas estadísticas o modelos predictivos que permitan evaluar la influencia conjunta de estos factores sobre el riesgo crediticio.